problem:  
Let \( (M^4, g) \) be a compact, simply connected Riemannian 4‑manifold with nonnegative sectional curvature, admitting an effective isometric circle action \( S^1 \curvearrowright M \).  The orbit space \(M^* = M/S^1\) is a smooth 3‑manifold with boundary, endowed with an Alexandrov metric of curvature \( \ge 0 \).  Denote by \( \Sigma^* \subset M^* \) the graph of singular orbits, consisting of edges corresponding to exceptional orbit 2‑spheres and vertices corresponding to fixed-point components (typically 2‑spheres).

By general theorems of Perelman–Petrunin, the orbit space of any isometric compact Lie group action is an **Alexandrov space** and the images of fixed-point sets are **extremal subsets**.  However, these results do not determine whether **1‑dimensional components of the singular graph**—arcs of exceptional orbits connecting two fixed-point components—are themselves *locally extremal* (i.e. totally geodesic in the quotient metric).  In all currently known examples (such as \( S^4, S^2 \times S^2, \mathbb{C}P^2, \mathbb{C}P^2 \# \pm\mathbb{C}P^2 \)), each edge of \( \Sigma^* \) projects to a genuine Alexandrov geodesic between boundary components.

**Problem:**  
Let \( \gamma^* \subset M^* \) be such an edge connecting boundary components of \(M^*\).  
> Must \( \gamma^* \) always be a local Alexandrov geodesic (i.e., an extremal subset of dimension one) in the quotient metric?

Equivalently:  
> Can nonnegative sectional curvature admit a smooth circle action whose orbit space contains an edge corresponding to exceptional orbits that is *not* an Alexandrov geodesic between its endpoint fixed components?

Why is it a "good" problem:  
This problem targets a subtle and currently unresolved interface between **Riemannian symmetry** and **Alexandrov extremal geometry**.  Known theorems (Perelman–Petrunin, Lytchak) ensure that images of fixed-point sets are extremal, but they leave open the extremality of *connecting 1‑dimensional strata* arising from orbits of intermediate isotropy.  

It is a “good” problem because:
- It isolates a concrete and elementary geometric phenomenon—geodesicity of singular arcs in the quotient—whose truth or failure would clarify how curvature on \(M\) rigidly shapes the quotient metric.  
- A positive answer would establish a new form of **metric rigidity** for nonnegatively curved manifolds with symmetry: all singular strata would be extremal, refining orbit‑space geometry.  
- A counterexample would reveal **unexpected flexibility** of Alexandrov orbit metrics compatible with nonnegative curvature, potentially generating new 4‑manifolds with exotic quotient geometry.  
- The question lies exactly at the frontier of **Alexandrov comparison geometry**, **transformation group theory**, and **curvature rigidity**, touching on how horizontal geodesics project and interact with isotropy strata.  
- The statement is compact and conceptually transparent—“Are exceptional orbit arcs always Alexandrov geodesics?”—but its resolution likely requires new synthesis of techniques from Alexandrov geometry and Riemannian group actions.

Thus the problem meets the “narrow entrance, profound depth’’ standard: an accessible geometric question whose answer could reshape the understanding of metric and curvature constraints in circle‑symmetric 4‑manifolds.